% Figure: Otto and Diesel cycles on the p-v plane, side by side. Otto has
% constant-volume heat addition (vertical 2->3); Diesel has constant-pressure
% heat addition (horizontal 2->3).
\begin{figure}[htbp]
\centering
\begin{tikzpicture}
% Otto (left)
\begin{axis}[
  name=otto,
  width=7.5cm, height=6.5cm,
  xlabel={specific volume $v$}, ylabel={pressure $p$},
  xmin=0, xmax=10, ymin=0, ymax=70,
  axis lines=left, xtick=\empty, ytick=\empty,
  label style={font=\small}, clip=false,
]
% r = 9: v2 = 1, v1 = 9
% 1->2 isentropic compression p1=1 at v=9 to p2 ~ 21.7 at v=1
\addplot[engblue, very thick, smooth] coordinates {(9,1) (5,2.5) (3,5.5) (1.8,11) (1,21.7)};
% 2->3 const-volume heat addition (vertical) v=1
\addplot[engred, very thick] coordinates {(1,21.7) (1,62)};
% 3->4 isentropic expansion v=1 to v=9
\addplot[engblue, very thick, smooth] coordinates {(1,62) (1.8,31) (3,15.5) (5,7) (9,2.86)};
% 4->1 const-volume heat rejection v=9
\addplot[engblue, very thick] coordinates {(9,2.86) (9,1)};
\filldraw[engorange] (axis cs:9,1) circle (1.4pt) node[anchor=north, font=\scriptsize]{1};
\filldraw[engorange] (axis cs:1,21.7) circle (1.4pt) node[anchor=east, font=\scriptsize]{2};
\filldraw[engorange] (axis cs:1,62) circle (1.4pt) node[anchor=south east, font=\scriptsize]{3};
\filldraw[engorange] (axis cs:9,2.86) circle (1.4pt) node[anchor=south, font=\scriptsize]{4};
\node[font=\scriptsize, anchor=south] at (axis cs:5,60) {(a) Otto};
\end{axis}
% Diesel (right)
\begin{axis}[
  name=diesel, at={($(otto.east)+(1.4cm,0)$)}, anchor=west,
  width=7.5cm, height=6.5cm,
  xlabel={specific volume $v$}, ylabel={pressure $p$},
  xmin=0, xmax=20, ymin=0, ymax=70,
  axis lines=left, xtick=\empty, ytick=\empty,
  label style={font=\small}, clip=false,
]
% r = 18: v2 = 1, v1 = 18. cut-off ratio 2.5: v3 = 2.5
% 1->2 isentropic compression p1=1 at v=18 to p2 ~ 57 at v=1
\addplot[engblue, very thick, smooth] coordinates {(18,1) (10,2.3) (5,6.0) (2.5,16.5) (1,57)};
% 2->3 const-pressure heat addition (horizontal) p=57, v=1 to v=2.5
\addplot[engred, very thick] coordinates {(1,57) (2.5,57)};
% 3->4 isentropic expansion v=2.5 to v=18
\addplot[engblue, very thick, smooth] coordinates {(2.5,57) (5,21) (10,8) (18,3.4)};
% 4->1 const-volume heat rejection v=18
\addplot[engblue, very thick] coordinates {(18,3.4) (18,1)};
\filldraw[engorange] (axis cs:18,1) circle (1.4pt) node[anchor=north, font=\scriptsize]{1};
\filldraw[engorange] (axis cs:1,57) circle (1.4pt) node[anchor=east, font=\scriptsize]{2};
\filldraw[engorange] (axis cs:2.5,57) circle (1.4pt) node[anchor=south west, font=\scriptsize]{3};
\filldraw[engorange] (axis cs:18,3.4) circle (1.4pt) node[anchor=south, font=\scriptsize]{4};
\node[font=\scriptsize, anchor=south] at (axis cs:10,60) {(b) Diesel};
\end{axis}
\end{tikzpicture}
\caption[Otto and Diesel cycles on the pressure-volume plane]{The two
reciprocating-engine cycles on the $p$-$v$ plane. (a) The Otto cycle adds heat
at constant volume (the vertical 2 to 3, modelling near-instantaneous
spark-ignited combustion at top dead centre) with compression ratio $r = 9$.
(b) The Diesel cycle adds heat at constant pressure (the horizontal 2 to 3,
modelling fuel injected and burned as the piston begins its downstroke) with
compression ratio $r = 18$ and cut-off ratio $r_c = 2.5$. Both reject heat at
constant volume (4 to 1). The Diesel tolerates a far higher compression ratio
because it compresses air alone, with no fuel present to pre-ignite.}
\label{fig:vol04ch03:otto-diesel-pv}
\end{figure}
